\section*{Capítulo \ref{ch:g11:color-light-sources} --- Cor e fontes de luz}

\begin{solution}{exo:g11:color-light-sources:1}
(a) amarelo; (b) ciano; (c) branco. Complementar do azul: o amarelo
(azul $+$ amarelo $=$ azul $+$ vermelho $+$ verde $=$ branco).
\end{solution}

\begin{solution}{exo:g11:color-light-sources:2}
Amarelo: V, Ve. Magenta: V, A. Branco: V, Ve, A. Preto: nenhum. Laranja:
V a plena \omterm{def:g10:energy-conservation:power}{potência}, Ve com \omterm{def:g10:energy-conservation:power}{potência} reduzida.
\end{solution}

\begin{solution}{exo:g11:color-light-sources:3}
O magenta absorve a faixa verde: sai $(V, A)$, luz magenta. O \omterm{def:g11:color-light-sources:subtractive}{filtro}
ciano absorve então a faixa vermelha: sobrevive só $(A)$ --- azul.
\end{solution}

\begin{solution}{exo:g11:color-light-sources:4}
Faixa amarela: \omterm{def:g11:color-light-sources:objectcolor}{difunde} V e Ve, absorve A. (a) Sob luz vermelha ela
\omterm{def:g11:color-light-sources:objectcolor}{difunde} V: vermelha. (b) Sob luz azul ela absorve tudo: preta.
\end{solution}

\begin{solution}{exo:g11:color-light-sources:5}
$T = \dfrac{\qty{2.90e-3}{m.K}}{\qty{5.80e-7}{m}} = \qty{5000}{K}$.
\end{solution}

\begin{solution}{exo:g11:color-light-sources:6}
\emph{(a)} $\lambda_{\max} = \qty{2.90e-3}{m.K}/\qty{1800}{K}
= \qty{1.61e-6}{m} = \qty{1610}{nm}$: \omterm{def:g10:light-spectra:wavelength}{infravermelho}.

\emph{(b)} O \omterm{def:g10:light-spectra:white}{espectro} é contínuo: a sua cauda visível é pequena, mas não
nula, e a ponta azul é a primeira a sumir --- o que resta de visível é
sobretudo vermelho-alaranjado, daí a cor da chama.
\end{solution}

\begin{solution}{exo:g11:color-light-sources:7}
Jaqueta (\omterm{def:g11:color-light-sources:objectcolor}{difunde} V): (a) vermelha; (b) vermelha; (c) preta. Luvas
(difundem Ve e A, absorvem V): (a) cianas; (b) pretas; (c) verdes.
\end{solution}

\begin{solution}{exo:g11:color-light-sources:8}
\emph{(a)} Um \omterm{def:g10:light-spectra:emission}{espectro de emissão} de linhas. Todo objeto só pode difundir
$\qty{589}{nm}$ ou nada: a rua fica em tons de amarelo-alaranjado e
preto.

\emph{(b)} O carro azul \omterm{def:g11:color-light-sources:objectcolor}{difunde} apenas a faixa azul. A luz do dia a
contém: carro azul. A linha do sódio não carrega azul algum: tudo é
absorvido, e o carro parece preto.
\end{solution}

\begin{solution}{exo:g11:color-light-sources:9}
\emph{(a)} O feixe branco se abre num arco-íris; o feixe do \omterm{def:g11:color-light-sources:lamps}{laser} é
desviado, mas continua um único feixe do mesmo verde: ele é
monocromático --- um só \omterm{def:g10:light-spectra:wavelength}{comprimento de onda}, e nada para o prisma
separar.

\emph{(b)} O \omterm{def:g11:color-light-sources:subtractive}{filtro} vermelho transmite apenas a faixa vermelha e absorve
$\qty{532}{nm}$: quase nada sai; o ponto desaparece.
\end{solution}

\begin{solution}{exo:g11:color-light-sources:10}
\omterm{def:g10:light-spectra:white}{Espectro}: um pico azul estreito perto de $\qty{450}{nm}$ mais a faixa
amarela larga do \omterm{def:g11:color-light-sources:led}{fósforo} (grosso modo, de $500$ a $\qty{700}{nm}$). Azul
e amarelo são \omterm{prop:g11:color-light-sources:complementary}{complementares}, de modo que o trio dos \omterm{def:g11:color-light-sources:cones}{cones} se lê como
branco. A faixa se apaga em direção ao vermelho intenso, e por isso um
objeto vermelho intenso recebe pouca coisa que possa difundir: ele parece
mais apagado do que à luz do dia, cujo \omterm{def:g10:light-spectra:continuous}{espectro contínuo} o alimenta por
inteiro.
\end{solution}

\begin{solution}{exo:g11:color-light-sources:11}
(a) $\lambda_{\max} = \qty{2.90e-3}{m.K}/\qty{1800}{K} = \qty{1610}{nm}$,
\omterm{def:g10:light-spectra:wavelength}{infravermelho}; (b) $\qty{2.90e-3}{m.K}/\qty{6500}{K} = \qty{446}{nm}$,
azul-violeta. A luz da vela é a chamada de ``quente'' --- a fonte de
temperatura \emph{mais baixa}: o vocabulário corre ao contrário da escala
\omterm{def:g10:light-spectra:kelvin}{kelvin}.
\end{solution}

\begin{solution}{exo:g11:color-light-sources:12}
\emph{(a)} \omterm{def:g11:color-light-sources:perception}{Metâmeras}. A existência delas prova que a cor é uma \omterm{def:g11:color-light-sources:perception}{percepção}
--- três respostas dos \omterm{def:g11:color-light-sources:cones}{cones} --- e não o \omterm{def:g10:light-spectra:white}{espectro} em si: infinitos
\omterm{def:g10:light-spectra:white}{espectros} levam ao mesmo trio.

\emph{(b)} Prisma ou \omterm{def:g10:light-spectra:white}{espectroscópio}: uma linha em $\qty{589}{nm}$ contra
duas faixas em $540$ e $\qty{610}{nm}$. \omterm{def:g11:color-light-sources:subtractive}{Filtro} vermelho: ele absorve
$\qty{589}{nm}$ e o ponto de sódio escurece, enquanto o subpixel de
$\qty{610}{nm}$ da tela sobrevive --- o trecho da tela fica vermelho.
\end{solution}

\begin{solution}{exo:g11:color-light-sources:13}
\emph{(a)} $\lambda_{\max} = \qty{2.90e-3}{m.K}/\qty{3200}{K} =
\qty{906}{nm}$. \omterm{def:g11:color-light-sources:colortemp}{Temperatura de cor} $\qty{3200}{K} < \qty{6500}{K}$: mais
quente (mais alaranjada) do que a luz do dia.

\emph{(b)} A gelatina absorve parte do excesso de vermelho-laranja para
reequilibrar o \omterm{def:g10:light-spectra:white}{espectro} em direção ao azul; como um \omterm{def:g11:color-light-sources:subtractive}{filtro} só retira luz,
a \omterm{def:g10:energy-conservation:power}{potência} total no palco cai.

\emph{(c)} Um \omterm{def:g11:color-light-sources:subtractive}{filtro} é passivo: ele pode absorver em cada \omterm{def:g10:light-spectra:wavelength}{comprimento de
onda}, mas nunca emitir, de modo que gelatina alguma consegue elevar a
intensidade do azul acima do que o refletor já produz.
\end{solution}

\begin{solution}{exo:g11:color-light-sources:14}
\emph{(a)} $\qty{2700}{K}$: luz quente, rica em vermelho --- a faixa
vermelha da carne é bem alimentada e a carne parece vividamente fresca.
$\qty{6500}{K}$: branco azulado e frio --- o peixe \omterm{def:g11:color-light-sources:objectcolor}{difunde} o trio
completo e parece branco e brilhante, recém-saído do gelo.

\emph{(b)} Nenhum dos dois mudou a mercadoria: os hábitos de absorção da
carne e do peixe continuam intactos. O que mudou foi a iluminação e,
portanto, a luz difundida --- a \omterm{def:g11:color-light-sources:objectcolor}{cor de um objeto} só se define sob uma
dada iluminação.
\end{solution}

\begin{solution}{exo:g11:color-light-sources:15}
\emph{(a)} O sódio: todo o brilho da cidade fica numa única linha em
$\qty{589}{nm}$, que um único \omterm{def:g11:color-light-sources:subtractive}{filtro} estreito retira das imagens do
telescópio. Um \omterm{def:g10:light-spectra:white}{espectro} largo de LED não pode ser filtrado sem descartar
junto a luz das estrelas.

\emph{(b)} Motoristas e pedestres preferem o LED: sob o seu \omterm{def:g10:light-spectra:white}{espectro}
completo os objetos mantêm cores distinguíveis (um casaco vermelho
continua vermelho), ao passo que a luz de sódio deixa tudo
amarelo-alaranjado ou preto.

\emph{(c)} LEDs de branco quente, em torno de $\qty{2700}{K}$: reprodução
de cores aceitável para a rua e \omterm{def:g10:light-spectra:white}{espectro} pobre na ponta azul, que é a
parte que mais polui o céu noturno.
\end{solution}

\begin{solution}{pb:g11:color-light-sources:1}
\textbf{1.} V $+$ Ve $=$ amarelo, V $+$ A $=$ magenta, Ve $+$ A $=$
ciano; a sobreposição das três é branca.

\textbf{2.} Laranja: vermelho a plena \omterm{def:g10:energy-conservation:power}{potência}, verde a \omterm{def:g10:energy-conservation:power}{potência} parcial,
azul apagado. Rosa-claro: os três acesos (base branca) com o vermelho
levemente dominante.

\textbf{3.} A sombra projetada em relação ao refletor vermelho recebe
apenas o refletor verde: sombra verde; simetricamente, a outra sombra é
vermelha.

\textbf{4.} A tinta azul-escura \omterm{def:g11:color-light-sources:objectcolor}{difunde} só A. Sob o refletor vermelho
sozinho: preta. Sob os três: ela \omterm{def:g11:color-light-sources:objectcolor}{difunde} a faixa azul --- azul.

\textbf{5.} O olho reduz qualquer \omterm{def:g10:light-spectra:white}{espectro} a três respostas dos \omterm{def:g11:color-light-sources:cones}{cones};
igualar o trio já basta para igualar a cor. Os três tipos de \omterm{def:g11:color-light-sources:cones}{cone} fixam o
número de refletores.

\textbf{6.} Amarelo $=$ branco $-$ azul: o corante absorve a faixa azul.

\textbf{7.} Refletor vermelho: \omterm{def:g11:color-light-sources:objectcolor}{difunde} V --- vermelho. Refletor verde:
\omterm{def:g11:color-light-sources:objectcolor}{difunde} Ve --- verde. Refletor azul: absorve tudo --- preto.

\textbf{8.} Ilumine a cena só com o refletor azul; o toque é de branco
$\to$ azul. O figurino amarelo \omterm{def:g11:color-light-sources:objectcolor}{difunde} V e Ve, mas absorve A: brilhante
sob o branco, preto sob o azul.

\textbf{9.} Gelatina magenta: $(V, Ve, A) \to (V, A)$. O figurino \omterm{def:g11:color-light-sources:objectcolor}{difunde}
V e Ve do que recebe: ele \omterm{def:g11:color-light-sources:objectcolor}{difunde} V e absorve A --- figurino vermelho.

\textbf{10.} Uma gelatina vermelha conserva uma faixa de três e joga fora
dois terços da \omterm{def:g10:energy-conservation:power}{potência} do refletor branco; já cada gelatina C, M ou Am
retira uma única faixa, conserva dois terços, e empilhar duas delas ainda
assim entrega qualquer primária.

\textbf{11.} Um \omterm{def:g10:light-spectra:emission}{espectro de emissão} de linhas. A rua vira um mundo
monocromático: toda superfície fica amarelo-alaranjada, mais apagada ou
mais viva, ou preta.

\textbf{12.} Parede branca: amarelo-alaranjada (\omterm{def:g11:color-light-sources:objectcolor}{difunde} a linha). Faixa
amarela: amarelo-alaranjada e viva. Carro vermelho: escuro --- a linha de
$\qty{589}{nm}$ fica fora da faixa que ele \omterm{def:g11:color-light-sources:objectcolor}{difunde}. Carro azul: preto.

\textbf{13.} Um único \omterm{def:g11:color-light-sources:subtractive}{filtro} estreito centrado em $\qty{589}{nm}$ retira
de uma imagem todo o brilho da cidade sacrificando quase nada do \omterm{def:g10:light-spectra:white}{espectro}
largo das estrelas.

\textbf{14.} $\lambda_{\max} = \qty{2.90e-3}{m.K}/\qty{2700}{K} =
\qty{1.07e-6}{m} = \qty{1070}{nm}$: o pico, e a maior parte da \omterm{def:g10:energy-conservation:power}{potência},
fica no \omterm{def:g10:light-spectra:wavelength}{infravermelho} --- calor, não luz.

\textbf{15.} Sódio: $0.30 \times \qty{100}{W} = \qty{30}{W}$ de luz
visível; tungstênio: $0.05 \times \qty{100}{W} = \qty{5}{W}$. Seis vezes
mais luz pela mesma conta.

\textbf{16.} O \omterm{def:g11:color-light-sources:led}{fósforo} converte parte do azul numa faixa amarela larga;
azul e amarelo são \omterm{prop:g11:color-light-sources:complementary}{complementares}, de modo que azul $+$ amarelo completa
o trio: branco.

\textbf{17.} Um pico azul estreito mais uma corcova amarela larga ---
contra o contínuo uniforme da luz do dia. A ponta do vermelho intenso (e
uma depressão entre o azul e o amarelo) fica sub-representada.

\textbf{18.} Os objetos de vermelho intenso saem apagados e
amarronzados. Daí os LEDs quentes de $\qty{2700}{K}$ do açougueiro, cujo
conteúdo reforçado de vermelho mantém a carne viva.

\textbf{19.} Nada lá dentro está a $\qty{2700}{K}$: o chip funciona perto
da temperatura ambiente. O rótulo promete apenas o \emph{tom} de um corpo
incandescente a $\qty{2700}{K}$, cujo pico ficaria em
$\qty{2.90e-3}{m.K}/\qty{2700}{K} = \qty{1.07e-6}{m} = \qty{1070}{nm}$.

\textbf{20.} Quarto: cerca de $\qty{2700}{K}$ --- luz quente e pobre em
azul, para descansar. Joalheria: $5500$ a $\qty{6500}{K}$ --- branco
parecido com a luz do dia, para que os brancos se leiam brancos e as
pedras devolvam todas as faixas. A lei: cor percebida $=$ \omterm{def:g10:light-spectra:white}{espectro} da
fonte $\times$ absorção do objeto, lido através de três \omterm{def:g11:color-light-sources:cones}{cones}.
\end{solution}
